\documentclass[10pt]{amsart}

\usepackage[margin=1in]{geometry}
\usepackage[T1]{fontenc}
\usepackage{lmodern}
\usepackage{amsmath,amssymb,amsthm,mathtools}
\usepackage{microtype}
\usepackage{url}
\usepackage{enumitem}
\usepackage[dvipsnames]{xcolor}
\usepackage[colorlinks=true,linkcolor=MidnightBlue,citecolor=OliveGreen,urlcolor=MidnightBlue]{hyperref}
\usepackage[capitalize,nameinlink]{cleveref}

\setlist{leftmargin=*,itemsep=2pt,topsep=4pt}
\emergencystretch=2.5em

\theoremstyle{plain}
\newtheorem{theorem}{Theorem}[section]
\newtheorem{lemma}[theorem]{Lemma}
\newtheorem{proposition}[theorem]{Proposition}
\newtheorem{corollary}[theorem]{Corollary}
\newtheorem{importedtheorem}[theorem]{Imported Theorem}
\newtheorem*{hypA}{Hypothesis $(\mathrm{H}_{\mathrm{A}})$}
\newtheorem*{hypO}{Hypothesis $(\mathrm{H}_{\mathrm{O}})$}
\theoremstyle{definition}
\newtheorem{definition}[theorem]{Definition}
\theoremstyle{remark}
\newtheorem{remark}[theorem]{Remark}
\crefname{importedtheorem}{Imported Theorem}{Imported Theorems}
\numberwithin{equation}{section}

\newcommand{\poly}{\operatorname{poly}}
\newcommand{\ENP}{\mathrm{E}^{\mathrm{NP}}}
\newcommand{\NTIME}{\mathrm{NTIME}}
\newcommand{\SIZE}{\mathrm{SIZE}}
\newcommand{\A}{\textup{[A]}}
\newcommand{\Oa}{\textup{[O]}}
\newcommand{\lean}[1]{\texttt{#1}}

\begin{document}

\title[Superlinear circuit lower bounds for $\mathrm{E}^{\mathrm{NP}}$]{Superlinear circuit lower bounds for $\mathrm{E}^{\mathrm{NP}}$,\\ assuming the new $\#$SAT algorithms}

\author{Vaibhav Tulsyan}

\date{November 2026}

\subjclass[2020]{68Q17, 68Q06, 68Q25}
\keywords{Circuit lower bounds, Exponential Time Hypothesis, algorithmic method, $\mathrm{E}^{\mathrm{NP}}$, satisfiability algorithms}

\begin{abstract}
In October 2026 OpenAI and Anthropic claimed deterministic subexponential-time algorithms for
$\#k$-SAT, refuting the Exponential Time Hypothesis. Anthropic's manuscript notes that, by a
theorem of Williams, this implies that $\ENP$ does not have circuits of linear size. We record the
quantitative form. Assuming the clause-count bound $2^{C\cdot kM/\sqrt{\log M}}$ of Anthropic's
Main Theorem, in the form of a Turing machine with engine constant $C$, there are $\varepsilon>0$
and a language $L_0\in\ENP$ such that for infinitely many input lengths $n$, $L_0$ has no circuit
of size $\varepsilon\,n\sqrt{\log n}$ over the full binary basis. The bound proved in OpenAI's Lean
development gives $\varepsilon\,n\sqrt{\log n}/(\log\log n)^2$ in the same way. The proof is
Williams' argument with the constants tracked: guess a witness circuit, check it with a circuit of
$3s+O(n)$ gates, apply Tseitin's transformation to get at most $4S+1$ clauses, and decide
satisfiability with the new algorithm; any $\varepsilon<1/(36aC)$ works, where $a$ is the
input-length shift of the construction. We explain why these algorithms give nothing at $n\log n$,
nothing for $\mathrm{NEXP}$ or $\mathrm{NQP}$, and nothing more from counting. The constant
$\varepsilon$ is not explicit, the bound holds only infinitely often, and everything is
conditional on the claimed refutations, which are under review.
\end{abstract}

\maketitle

% =====================================================================
\section{Introduction}\label{sec:intro}
% =====================================================================

Two manuscripts, by OpenAI~\cite{OpenAI26} and by Anthropic~\cite{Anthropic26} (below \Oa{} and
\A), each claim a deterministic algorithm that counts the satisfying assignments of a $k$-CNF in
subexponential time, and hence a refutation of the Exponential Time Hypothesis (ETH). Both come
with Lean~4 formalisations of their top-level statements. These are \emph{claimed} results, still
under review. The Main Theorem of \A{} (Theorem~1, made explicit in its Corollary~11) counts the
models of a $k$-CNF with $m\ge2$ clauses in which $n\ge4$ variables occur in time
\[
  \min\bigl\{2^{C\cdot km/\sqrt{\log m}},\;2^{C\cdot k^2 n\log\log n/\sqrt{\log n}}\bigr\}
\]
for an absolute constant $C$. Logarithms are to base~$2$.

Williams~\cite[Theorem~6.1]{Wil13} showed that if ETH is false, then $\ENP$ does not have circuits
of linear size; see also Jahanjou, Miles and Viola~\cite[Corollary~2(1) with Corollary~3(1$'$)]{JMV15}.
\A{} draws this consequence and formalises it in Lean (\lean{ENP.CorCircuits}). Its proof uses only
the qualitative statement that, for every $\delta>0$, $3$-SAT is decidable in time $2^{\delta N}$.
It gives ``not $O(n)$'' with no rate. Here we use the rate.

\subsection{Model}
A \emph{circuit} is over the full binary basis $B_2$: each gate computes one of the $16$
functions $\{0,1\}^2\to\{0,1\}$ of two earlier wires. Its \emph{size} is the number of gates;
inputs are not counted. Circuits are non-uniform, one per input length. For $s\colon\mathbb N\to\mathbb N$,
a language $L$ \emph{has circuits of size $s$ at length $n$} if some circuit with at most $s(n)$
gates computes $L\cap\{0,1\}^n$. $\ENP$ is the class of languages decided in time $2^{O(n)}$ by a
deterministic oracle Turing machine with an oracle for $\mathrm{SAT}$, where queries may have
length $2^{O(n)}$.

\begin{hypA}
There are a constant $C\ge1$ and a deterministic multitape Turing machine that, given a $3$-CNF
$\psi$ with no repeated clause, with $M\ge2$ clauses and at least $4$ occurring variables, encoded
in length $L$, decides whether $\psi$ is satisfiable within $L^{C}\cdot2^{3C\cdot M/\sqrt{\log M}}$
steps.
\end{hypA}

\noindent This is the $m$-branch of Corollary~11 of \A with $k=3$, stated for Turing machines. The
constant $C$ here is the \emph{engine constant}, the constant of the machine that is actually run.
It is not necessarily the $C$ printed in \A. \A{} counts operations in its own model, and turning
those into Turing-machine steps changes $C$ by the polynomial overhead of the simulation. In
Anthropic's Lean development the rate is available as the machine form \lean{Tied.TopM}
(rate $e_M\le G\cdot k\cdot m/\sqrt{\log(m+2)}$, then a polynomial-time machine on a pad of length
$2^{e_M}$). The engine constant is then determined by $G$ and by the polynomial exponent of that
machine.

\begin{hypO}
There are a constant $C_2\ge1$ and a deterministic multitape Turing machine that decides the
satisfiability of a $3$-CNF with $N\ge 16$ variables, encoded in length $L$, within
$L^{C_2}2^{C_2\rho(N)}$ steps, where $\rho(N)=N(\log\log N)^2/\sqrt{\log N}$.
\end{hypO}

\noindent This is the form of the theorem \lean{exact\_count\_and\_sat} in OpenAI's Lean
development, whose rate function is $\rho$, up to the base of the logarithms. The printed main
theorem of \Oa{} has the exponent $O_k(n\log\log n/\sqrt{\log n})$, with one factor of
$\log\log n$. It would give the bound of \cref{cor:openai} with $(\log\log n)^2$ replaced by
$\log\log n$, in the same way.

\begin{theorem}[Main theorem; written proof, not machine-checked]\label{thm:main}
Assume $(\mathrm{H}_{\mathrm{A}})$. There are $\varepsilon>0$ and $L_0\in\ENP$ such that for every
$N$ there is an $n\ge N$ for which $L_0\cap\{0,1\}^n$ is computed by no circuit over $B_2$ with at
most $\varepsilon\,n\sqrt{\log n}$ gates.

More precisely, there is a constant $a\ge2$, which depends only on the succinct reduction of
\cref{imp:jmv} applied to one fixed language, such that every $\varepsilon$ with
$36\,a\,C\,\varepsilon<1$ works.
\end{theorem}

\begin{corollary}[Through OpenAI's bound]\label{cor:openai}
Assume $(\mathrm{H}_{\mathrm{O}})$. There are $\varepsilon>0$ and $L_0\in\ENP$ such that for
infinitely many $n$, $L_0\cap\{0,1\}^n$ has no circuit with at most
$\varepsilon\,n\sqrt{\log n}/(\log\log n)^2$ gates. Any $\varepsilon$ with
$3aC_2\varepsilon<1$ works.
\end{corollary}

The $n$-branch of \A{} is a bound in the number of variables, $2^{C k^2N\log\log N/\sqrt{\log N}}$.
Used in the same way, it gives only $\varepsilon\,n\sqrt{\log n}/\log\log n$. So we use the
$m$-branch.

\begin{remark}[What the statement does and does not say]\label{rem:form}
\leavevmode
\begin{enumerate}[label=\textup{(\alph*)}]
\item \emph{Infinitely often.} The conclusion is ``$\forall N\,\exists n\ge N$''. It is not that
$L_0$ fails at a single length. With $\varepsilon$ small, $\lfloor\varepsilon n\sqrt{\log n}\rfloor=0$
for many small $n$, so a single failing length would be close to vacuous. For a fixed $\varepsilon$,
the infinitely-often form cannot be recovered from the single-length form by enlarging the size
bound, which is what the linear case does (\lean{ENP.infinitely\_often}). We prove it directly by
assuming circuits only at lengths $\ge N$. The conclusion does not say that the bound holds for
almost every $n$ (\cref{sec:limits}).
\item \emph{An existential $\varepsilon$.} The constant $\varepsilon$ is not explicit, because $C$
is not: \A{} and its Lean statement only assert that some $C$ exists. The result therefore does not
give $c\,n\sqrt{\log n}$ for every $c$. That would need a rate $2^{o(M/\sqrt{\log M})}$.
\item \emph{Strength.} The theorem implies the linear-size statement of \A, and it is strictly more
informative: at infinitely many lengths, the lower bound grows faster than every linear function.
\end{enumerate}
\end{remark}

\subsection{What is new, and what is not}
The qualitative consequence ``$\ENP\not\subseteq\SIZE(O(n))$'' is stated, and formalised, in \A.
JMV note that their proofs ``provide a close correspondence between the running time for
satisfiability and the parameters of the circuit class''~\cite{JMV15}. Ben-Sasson and Viola's
constant-level instance, as quoted there, is: $3$-SAT in time $c^n$ with $c<2^{1/10}$ implies
$\ENP\not\subseteq\SIZE(3n)$. What we add is the substitution of the rate
$2^{O(M/\sqrt{\log M})}$, the bookkeeping that keeps every overhead linear, and the resulting
explicit superlinear rate. We have not found this rate stated elsewhere. Given how direct it is, it
may well have been noticed independently. We claim the careful statement and proof, not the idea.

% =====================================================================
\section{Imported results}\label{sec:imported}
% =====================================================================

\begin{importedtheorem}[Nondeterministic time hierarchy, \v{Z}\'ak~\cite{Zak83}]\label{imp:hier}
If $t_2$ is time-constructible and $t_1(n+1)=o(t_2(n))$, then
$\NTIME[t_1]\subsetneq\NTIME[t_2]$. In particular, for every $\gamma<1$ there is a language
$L^\star\in\NTIME[2^n]\setminus\NTIME[2^{\gamma n}]$.
\end{importedtheorem}

\begin{importedtheorem}[Local reductions, {\cite[Theorem~1 and the remark after it]{JMV15}}]\label{imp:jmv}
Let $L\in\NTIME[2^n]$. There are constants $d,a_1$ and a polynomial-time map sending
$x\in\{0,1\}^n$ to a circuit $E_x\colon\{0,1\}^{u}\to\{0,1\}^{3v+3}$, with $u=n+O(\log n)$ and
$v\le a_1 n$, such that the following hold.
\begin{itemize}
\item $E_x(j)$ lists the three variables (each in $\{0,1\}^v$) and three signs of the $j$-th clause
of a $3$-CNF $\varphi_x$ with $2^u$ clauses on the variable set $\{0,1\}^v$.
\item $x\in L$ if and only if $\varphi_x$ is satisfiable.
\item Every output bit of $E_x$ is a decision tree of depth $O(1)$ in the bits of $j$. Hence $E_x$
has at most $d\,n$ gates over $B_2$.
\end{itemize}
\end{importedtheorem}

\noindent JMV state this for algorithms $M(x,y)$ running in time $T\ge2^{\Omega(n)}$. Their index
length is $\ell=\log T+O(\log n)$, and every output bit is a decision tree of depth $O(\log r)$
with $r=n/\log T=O(1)$. Williams' \cite[Theorem~6.1]{Wil13} is a \emph{sketch} with the same
shape. Its formula comes from a tableau of size $O(2^{3n})$, and its check circuit has ``$O(cn)$
wires and $3n+O(1)$ inputs''. The explicit count $3s+O(n)$ that we use is the one in JMV's \S6
and in Anthropic's Lean module (\lean{StForm}). There the reduction is replaced by a
self-contained succinct tableau formula and lazy diagonalisation. We may pad the variable names so
that $v=v(n)$ depends only on $n$ and $n\mapsto n+v(n)$ is strictly increasing. Put
\[
  \ell(n):=n+v(n),\qquad a:=\max\{2,\;1+a_1\},\qquad\text{so}\qquad \ell(n)\le a n.
\]

% =====================================================================
\section{Proof of \texorpdfstring{\cref{thm:main}}{Theorem 1.1}}\label{sec:proof}
% =====================================================================

We follow Williams' route~\cite[Lemma~3.2, Theorem~6.1]{Wil13}. Fix $\varepsilon>0$ with
$36aC\varepsilon<1$, and fix $\gamma$ with $36aC\varepsilon<\gamma<1$. Let $L^\star$ be the
language of \cref{imp:hier} for this $\gamma$, and let $E_x,\varphi_x,u,v,\ell$ be as in
\cref{imp:jmv} for $L=L^\star$. Put
\[
  s(\ell):=\bigl\lfloor\varepsilon\,\ell\sqrt{\log\ell}\bigr\rfloor .
\]

\subsection{The language \texorpdfstring{$L_0$}{L0}}
For $x$ with $\varphi_x$ satisfiable, let $w^\ast(x)\in\{0,1\}^{2^{v}}$ be its lexicographically
first satisfying assignment, indexed by $\{0,1\}^{v}$. Let $L_0$ be the set of strings $xi$ with
$|x|=n$, $i\in\{0,1\}^{v(n)}$, $\varphi_x$ satisfiable, and $w^\ast(x)_i=1$. Strings whose length is
not of the form $\ell(n)$ are not in $L_0$. Since $\ell$ is strictly increasing, $n$ and the
split $xi$ can be read off from the length.

\begin{lemma}\label{lem:enp}
$L_0\in\ENP$.
\end{lemma}

\begin{proof}
This is \cite[Lemma~3.2]{Wil13}. On input $xi$ of length $\ell(n)$, write out $\varphi_x$ by
evaluating $E_x$ on all $2^u$ indices, which takes time $2^{O(n)}$. The language
\[
  \{\langle\varphi,p\rangle:\varphi\text{ has a satisfying assignment extending the partial assignment }p\}
\]
is in $\mathrm{NP}$. Use it to fix the bits of $w^\ast(x)$ one at a time, in
lexicographic order. This takes $2^{v}$ queries of length $2^{O(n)}$, after one query that rejects
if $\varphi_x$ is unsatisfiable. Then output bit $i$. The total time is $2^{O(n)}=2^{O(\ell)}$.
\end{proof}

\subsection{The algorithm}
Suppose, for a contradiction, that there is an $N$ such that $L_0$ has circuits of size $s$ at every
length $\ell\ge N$. Let $n$ be so large that $\ell(n)\ge N$, and let $x\in\{0,1\}^n$. Consider the
following nondeterministic algorithm.

\begin{enumerate}[label=\textup{(\arabic*)}]
\item \emph{Guess} a circuit $K$ on $\ell(n)$ inputs with at most $s(\ell(n))$ gates. Hardwire $x$
to get $W(i):=K(xi)$, a circuit on $v$ inputs with at most $s(\ell(n))$ gates (substituting
constants does not add gates).
\item \emph{Build the check circuit} $D\colon\{0,1\}^u\to\{0,1\}$. On input $j$, compute
$E_x(j)=(y_1,\sigma_1,y_2,\sigma_2,y_3,\sigma_3)$. Evaluate the three copies $W(y_1),W(y_2),W(y_3)$.
Compute the three literal values $W(y_t)\oplus\neg\sigma_t$ (one gate each), and output $1$ if all
three are $0$ (two gates). So $D(j)=1$ if and only if clause $j$ of $\varphi_x$ is violated by the
assignment $i\mapsto W(i)$. Its size is
\begin{equation}\label{eq:S}
  S\;\le\;3\,s(\ell(n))+d\,n+5 .
\end{equation}
\item \emph{Tseitin.} Introduce one variable per input and per gate of $D$. For a gate
$g=f(p,q)$, add for each $(\alpha,\beta)\in\{0,1\}^2$ the clause
$(p\ne\alpha)\vee(q\ne\beta)\vee(g=f(\alpha,\beta))$, of width at most $3$. Add the unit clause on
the output gate. Delete tautological and repeated clauses. The result $\psi$ is a $3$-CNF with no
repeated clause and
\begin{equation}\label{eq:M}
  M\;\le\;4S+1
\end{equation}
clauses. It is satisfiable if and only if $D(j)=1$ for some $j$, since the gate variables are
forced. The number of variables is $u+S$, so also $M\le4N_v+1$, where $N_v$ is the number of
variables.
\item \emph{Decide.} If $M\le1$ or fewer than $4$ variables occur, decide $\psi$ by brute force.
Otherwise run the machine of $(\mathrm{H}_{\mathrm{A}})$. \emph{Accept} if and only if $\psi$ is
unsatisfiable.
\end{enumerate}

\begin{lemma}[Correctness]\label{lem:correct}
For $|x|=n$ as above, the algorithm has an accepting branch if and only if $x\in L^\star$.
\end{lemma}

\begin{proof}
If $x\in L^\star$, then $\varphi_x$ is satisfiable. Some circuit $K$ of size $\le s(\ell(n))$
computes $L_0$ at length $\ell(n)$, so on the branch that guesses it,
$W(i)=w^\ast(x)_i$ for every $i$. That assignment violates no clause, so $D\equiv0$, $\psi$ is
unsatisfiable, and the branch accepts. If $x\notin L^\star$, then every assignment, in particular
$i\mapsto W(i)$ for every guessed $K$, violates some clause $j$. So $D(j)=1$, $\psi$ is satisfiable,
and every branch rejects.
\end{proof}

\subsection{The arithmetic}
\begin{lemma}[Running time]\label{lem:time}
For all sufficiently large $n$, every branch halts within $2^{\gamma n}$ steps.
\end{lemma}

\begin{proof}
Steps (1)--(3) take $\poly(n,s)=2^{O(\log n)}$ steps, and the encoding length of $\psi$ is
$L=\poly(n)$. It remains to bound the exponent $3C\,M/\sqrt{\log M}$ of step (4).

By \eqref{eq:S}, \eqref{eq:M} and $\ell(n)\le an$,
\[
  M\;\le\;M^\ast:=12\,\varepsilon\,a\,n\sqrt{\log(an)}+4dn+21 .
\]
The function $x\mapsto x/\sqrt{\log x}$ is increasing for $x\ge2$, since its derivative is
$(\log x)^{-1/2}\bigl(1-\tfrac{1}{2\ln x}\bigr)>0$ there. Also $M^\ast\ge n$, so
$\log M^\ast\ge\log n$. Hence
\begin{equation}\label{eq:exp}
  \frac{3C\,M}{\sqrt{\log M}}\;\le\;\frac{3C\,M^\ast}{\sqrt{\log n}}
  \;=\;36\,a\,C\,\varepsilon\,n\sqrt{\frac{\log(an)}{\log n}}+\frac{3C(4dn+21)}{\sqrt{\log n}}
  \;=\;36\,a\,C\,\varepsilon\,n+o(n).
\end{equation}
So the running time is at most $2^{O(\log n)}\cdot L^C\cdot2^{(36aC\varepsilon+o(1))n}$, which is
at most $2^{\gamma n}$ for large $n$, because $36aC\varepsilon<\gamma$.
\end{proof}

\begin{proof}[Proof of \cref{thm:main}]
The algorithm needs to know $s(\ell(n))$, which is computable in time $\poly(n)$. Under the
assumption, \cref{lem:correct,lem:time} decide $L^\star$ correctly within $2^{\gamma n}$ steps on all
inputs of all sufficiently large lengths. The finitely many remaining inputs are handled by a
table. Hence $L^\star\in\NTIME[2^{\gamma n}]$, which contradicts \cref{imp:hier}. So for every $N$
there is an $\ell\ge N$, indeed one of the form $\ell(n)$, at which $L_0$ has no circuit with at
most $s(\ell)$ gates. Together with \cref{lem:enp}, this is the theorem.
\end{proof}

\begin{remark}[The input-length shift and the constant]\label{rem:constant}
\leavevmode
\begin{enumerate}[label=\textup{(\alph*)}]
\item \emph{The shift.} The witness circuit is a circuit for $L_0$ at length $\ell(n)$, not $n$.
So the check circuit has $3s(\ell(n))+O(n)$ gates, not ``$3s(n)+O(n)$''. In JMV's own
construction the $\ENP$ input is $(x,i)$ of length $2n+O(\log n)$. For $s=\varepsilon
n\sqrt{\log n}$ the shift costs only the factor $a$ in the condition on $\varepsilon$.
\item \emph{Lean normalisation.} Anthropic's module avoids the hierarchy theorem. It takes $m$ to be
the input length of its $\ENP$ language. The tableau formula there has time $2^{b_t}$ with
$b_t\ge3m/8$, and the engine is allowed a budget of $2^{m/32}$ inside a total of $2^{m/8}$
(\lean{exists\_threshold}). Its check formula has $N_v\le3|K|+a_0(m+1)+g(e)$ variables and at most
$4N_v+1$ clauses. With $|K|\le\varepsilon(m+1)\sqrt{\log(m+2)}$, the computation
\eqref{eq:exp} becomes $36C\varepsilon m(1+o(1))+O(m/\sqrt{\log m})\le m/32-o(m)$. This holds for
$\varepsilon<1/(1152\,C)$; for instance $\varepsilon=1/(1200C)$.
\item \emph{Which $C$.} In either normalisation, $C$ is the effective engine constant of
$(\mathrm{H}_{\mathrm{A}})$. In Lean this is $G$ together with the polynomial exponent of the
machine on the pad, and possibly the exponent of the step bound in \lean{PaddedDescent.exists\_tm'}.
It is fixed before the formula table is chosen, so $\varepsilon$ is still an absolute constant.
\end{enumerate}
\end{remark}

\begin{proof}[Proof of \cref{cor:openai}]
Take $s(\ell)=\lfloor\varepsilon\ell\sqrt{\log\ell}/(\log\log\ell)^2\rfloor$ and $3aC_2\varepsilon<\gamma<1$.
Run the same algorithm with the machine of $(\mathrm{H}_{\mathrm{O}})$ in step (4). Now
$N_v\le u+S\le N^\ast:=3\varepsilon an\sqrt{\log(an)}/(\log\log n)^2+(d+2)n$ for large $n$. The
function $\rho$ is increasing for large arguments. We have $\log N^\ast=\log n+O(\log\log n)$, and
so $\log\log N^\ast=\log\log n+o(1)$. Therefore
\[
  C_2\,\rho(N_v)\le C_2\,\rho(N^\ast)=3aC_2\varepsilon\,n(1+o(1))+O\!\left(\frac{n(\log\log n)^2}{\sqrt{\log n}}\right)\le(\gamma-o(1))\,n,
\]
and the rest is unchanged.
\end{proof}

% =====================================================================
\section{Comparison with what is known}\label{sec:known}
% =====================================================================

The following statements were checked against the primary sources.
\begin{itemize}
\item \emph{$\ENP$ and $\mathrm{P}$ over $B_2$.} Before 2026 the best lower bound for a language in
$\ENP$ was $3.1n-o(n)$, by Li and Yang~\cite{LY22}. They prove it for affine dispersers of
sublinear dimension, which are constructible in polynomial time, so the bound holds already in
$\mathrm{P}$. It improves the $(3+\tfrac1{86})n-o(n)$ of Find, Golovnev, Hirsch and
Kulikov~\cite{FGHK16}, which in turn improves the $3n-o(n)$ of Blum. Williams wrote in 2010 that
``linear size circuit lower bounds are not known for this class''~\cite{Wil13}, that is, no proof
that linear size does not suffice.
\item \emph{Above $\ENP$.} Chen, Hirahara and Ren~\cite{CHR24} prove that
$\mathrm{S}_2\mathrm{E}/_1\not\subseteq\SIZE[2^n/n]$. This gives the same near-maximum bounds for
$\Sigma_2\mathrm{E}$, $(\Sigma_2\mathrm{E}\cap\Pi_2\mathrm{E})/_1$ and $\mathrm{ZPE}^{\mathrm{NP}}/_1$.
The almost-everywhere, advice-free bound $\mathrm{S}_2\mathrm{E}\not\subseteq\text{i.o.-}\SIZE[2^n/n]$
is due to Li~\cite{Li24}. $\ENP\subseteq\mathrm{S}_2\mathrm{E}$, so these classes lie above the
one treated here.
\item \emph{Fixed-polynomial bounds.} Kannan~\cite{Kannan82} showed that for each $k$,
$\Sigma_2\mathrm{P}\cap\Pi_2\mathrm{P}\not\subseteq\SIZE(n^k)$. Whether $\ENP\not\subseteq\SIZE(n^k)$
for some fixed $k>1$ is open. Before the refutations, so was $\ENP\not\subseteq\SIZE(O(n))$.
\item \emph{$\mathrm{NEXP}$ and linear size.} No superlinear lower bound for $\mathrm{NEXP}$ follows.
JMV raise as an open problem the easy-witness strengthening of~\cite{IKW02} to linear-size
circuits~\cite{JMV15}.
\end{itemize}
So, conditionally, \cref{thm:main} is a superlinear lower bound with an explicit rate for a class
below $\mathrm{S}_2\mathrm{E}$ over $B_2$, holding infinitely often. It is the first we know of,
but we cannot rule out that it has been stated elsewhere.

% =====================================================================
\section{What these algorithms cannot give}\label{sec:limits}
% =====================================================================

\begin{itemize}
\item \emph{Beyond $n\sqrt{\log n}$.} The check circuit and Tseitin's transformation are already
linear. So the only lever is the rate: an algorithm in time $2^{M/g(M)}$ converts into
$s=\varepsilon n\,g(n)$. Size $n\log n$ would need $3$-SAT in $2^{O(M/\log M)}$, and size
$n^{1+\delta}$ would need $2^{M^{1-\delta'}}$. Neither manuscript claims such a rate.
\item \emph{$\mathrm{NEXP}$ and $\mathrm{NQP}$.} Moving below $\ENP$ needs an easy-witness lemma.
Those of Impagliazzo, Kabanets and Wigderson~\cite{IKW02} and of Murray and Williams~\cite{MW18}
work for polynomial size. The second loses a polynomial in the size. After the loss, one needs
satisfiability of $\poly(n)$-size circuits in time $2^n/n^{\omega(1)}$, and $2^{O(s/\sqrt{\log s})}$
with $s=n^k$ gives nothing. Linear-size easy witnesses are JMV's open problem.
\item \emph{Counting adds nothing here.} The $\mathrm{NP}$ oracle of $\ENP$ is used only to fix the
lexicographically first witness, and the final check is a decision question. The counting form of
the algorithms is not used. The lower bounds of the algorithmic method that do exploit counting or fast
evaluation, such as those against $\mathrm{ACC}$ of Williams~\cite{Wil14} and
Murray--Williams~\cite{MW18}, concern polynomial-size circuits and run into the previous item.
\item \emph{Almost everywhere.} Chen, Lyu and Williams~\cite{CLW20} obtain almost-everywhere lower
bounds for $\ENP$ (against $\mathrm{ACC}^0$) from non-trivial algorithms. Whether their machinery
keeps the check circuit at $O(s)$ gates, so that it would give an almost-everywhere version of
\cref{thm:main}, we have not checked. Their reductions involve PCPs, whose polylogarithmic overhead
would absorb the $\sqrt{\log n}$.
\item \emph{Restricted classes.} The algorithms are CNF algorithms, and the reduction goes through
general circuits. The implied bounds for formulas or for log-depth circuits are no better than what
\cref{thm:main} gives, and for formulas they are weaker than known bounds.
\end{itemize}

% =====================================================================
\section{Formalisation status}\label{sec:formal}
% =====================================================================

Anthropic's Lean development~\cite{Anthropic26} proves the linear-size form,
\[
  \lean{ENP.CorCircuits}:\ \exists L,\ \lean{InENP}\ L\wedge\neg\,\lean{HasLinearCircuits}\ L,
\]
together with the infinitely-often variant \lean{CorCircuitsInf}. Its circuits are over $B_2$, and
size is the number of gates. $\ENP$ is formalised with Mathlib \lean{FinTM2} oracle machines.
The proof follows the route above, using the $\delta$-form of the engine (time $c\cdot2^{\delta N}$
for every $\delta>0$) and lazy diagonalisation in place of \cref{imp:hier}. The statement
\lean{HasCircuits}~$L$~$s$ is already parametric in $s$, but the module hard-wires both the size
$c_0(m+1)$ and the time $2^{\delta N}$.

A Lean formalisation of the $n\sqrt{\log n}$ form on top of that module is in progress.

\Cref{thm:main} and \cref{cor:openai} are therefore proved here on paper only. We do not claim that
they are machine-checked. The changes needed are recorded in the project
repository~\cite{TulsyanRepo26}. Three points are not purely arithmetic.
\begin{itemize}
\item The engine's states must be allowed to grow with the rate, not polynomially in the input:
the rate machine runs on a pad of length $2^{e_M}$.
\item The engine interface must be weakened to a decision specification, or its count fixed up
over non-occurring variables.
\item The target must be stated in the infinitely-often form of \cref{rem:form}(a).
\end{itemize}

% =====================================================================
\section{Limitations}\label{sec:limitations}
% =====================================================================

\begin{enumerate}[label=\textup{(\alph*)}]
\item \emph{Conditional.} Everything is conditional on $(\mathrm{H}_{\mathrm{A}})$, or on
$(\mathrm{H}_{\mathrm{O}})$ for \cref{cor:openai}. These are Turing-machine forms of the claimed
theorems of \A{} and \Oa. Expert review of the manuscripts is pending, and if the hypotheses fail,
the results lose their proofs. Our proofs are on paper, written with AI assistance, and not yet
checked by a human.
\item \emph{$\varepsilon$ is not explicit.} It depends on the effective engine constant $C$,
which is existential in \A{} and in its Lean statement, and on the constant $a$ of the reduction.
No numerical value of $\varepsilon$ is claimed. The bound gives no meaningful statement at any
concrete input length.
\item \emph{Infinitely often only.} The lower bound holds at infinitely many lengths, not at almost
every length (\cref{sec:limits}).
\item \emph{Imported reduction.} \Cref{imp:jmv} is used as stated by JMV, with the linear gate
count read off from depth-$O(1)$ decision trees. Williams' Theorem~6.1, which we follow, is itself
a sketch. The self-contained version is Anthropic's Lean module, for the linear case.
\item \emph{Literature.} The statements in \cref{sec:known} and all citations were checked against
the primary sources (publisher records, arXiv, ECCC, the author's version of~\cite{Wil13}) on
4~November~2026. We have not searched systematically for prior statements of \cref{thm:main}.
\end{enumerate}

\section*{Acknowledgements}
This work was done with AI assistance (Claude Code, Anthropic): the observation, the proofs, the
referee check of the first draft and the literature checks were produced in sessions with Claude
Code. Everything here has been checked against the primary sources as described, but awaits human
expert review. Responsibility for any errors rests with the author. This note continues the
project of~\cite{Tulsyan26}.

\bibliographystyle{amsplain}
\bibliography{refs}

\end{document}
